\begin{proposition}[E.27 --- Proved: the canonical rational splitting and the three Kummer coordinates]
\label{E-explicit-kernel}
Let $\mathcal G=G_{\mathbb Q,\{3,\infty\}}$, and retain
$\mathcal H$, $B_w$, $\mathcal A$, and $A^+$ from E.23--E.25.
The rational averaging operator
\[
 e_{\mathcal A}=\frac1{54}\sum_{a\in\mathcal A}a
       \quad\text{on }A_{\mathbb Q_3}
\]
gives the following specified direct-sum decomposition:
\[
 \mathcal H=
 \operatorname{im}H^1(\mathcal G,V^+_{\mathbb Q_3})
 \ \oplus\
 H^1\bigl(\mathcal G,(A_{\mathbb Q_3})^{\mathcal A}(-1)\bigr).
 \tag{E27.1}
\]
The projection to the second summand is induced by
$e_{\mathcal A}$. In the cusp basis of E.18 put
\[
 b_0=e_{A_0}+e_{B_0},\quad
 b_1=e_{A_1}+e_{B_1},\quad
 b_2=e_{A_2}+e_{B_2},\quad
 b_3=e_{C_1}+e_{C_2}.
\]
There is an integral decomposition with its indicated bases,
\[
 V^+=
 \mathbb Z_3(1)\{b_0,b_1+b_2,b_3\}
 \ \oplus\
 \mathbb Z_3(1)\varepsilon_K\{b_1-b_2\},
 \tag{E27.2}
\]
and
\[
 H^1(\mathcal G,V^+)=\mathbb Z_3^3\oplus\mathbb Z/3.
 \tag{E27.3}
\]
After rationalization the first three coordinates in
(E27.1) are therefore uniquely specified scalar multiples
of $\kappa_3(3)$ in the three displayed Tate directions,
mapped to $A(-1)$ by $v\mapsto[v,w]$.
The other two coordinates are the classes in the two
one-dimensional cohomology summands identified in E.23.
This proposition fixes their normalization without
evaluating the actual class. E.28 eliminates one
projected coordinate, E.29 relates two Tate coordinates,
and E.35--E.36 identify the elliptic class and evaluate
the remaining two rational Kummer coordinates.
\end{proposition}
\begin{proof}
The action of $\mathcal A$ on the graded module $A$ is
genuine and commutes with arithmetic up to the permutation
of elements of $\mathcal A$. Thus its average commutes
with $\mathcal G$. It preserves $B_w$, on which the
fixed subspace is
$V_{\mathbb Q_3}^{\mathcal A}(1)=0$ by E.23.
Applying this exact rational projector to
$0\to B_w\to A\to\mathcal Q\to0$ gives an isomorphism
\[
       (A_{\mathbb Q_3})^{\mathcal A}
                \xrightarrow{\sim}
       (\mathcal Q_{\mathbb Q_3})^{\mathcal A}=D.
\]
The isomorphism is the actual quotient map. In E.25,
the image of a class in $\mathcal H$ lies in $H^1(D(-1))$.
Subtracting its $e_{\mathcal A}$-projection therefore
places it in the kernel of the quotient map on $H^1$.
That kernel is precisely the injected image of
$H^1(V^+_{\mathbb Q_3})$ by the exact sequence of E.25.
The two summands have zero intersection, since
$e_{\mathcal A}$ is zero on the first and the identity
on the second. This proves (E27.1) with its projection.

For (E27.2), the compact coordinates of an
$I_V$-fixed vector satisfy $-h=h$ and hence vanish.
Its toric coordinates are exactly the four paired
cusp vectors $b_i$ above. Every cusp is rational over
$K=\mathbb Q(\zeta_3)$, so $G_K$ acts on them by
$\chi$. Complex conjugation satisfies
\[
 cb_0=-b_0,\quad cb_3=-b_3,\quad
 cb_1=-b_2,\quad cb_2=-b_1.
\]
The basis change from $(b_1,b_2)$ to
$(b_1+b_2,b_1-b_2)$ has determinant $-2$, a
$3$-adic unit. The first three directions have the
ordinary Tate action, and the last has its quadratic
twist, proving the integral assertion.

Kummer theory over $\mathbb Z[1/3]$ gives
$H^1(\mathcal G,\mathbb Z_3(1))=\mathbb Z_3\kappa_3(3)$:
the units are $\{\pm3^n\}$ and the class group is zero.
For the twisted line restrict to $G_{K,\{3\}}$ and use
the exact order-two averaging projector. The ring
$\mathbb Z[\zeta_3]$ is Euclidean for its norm; its
units after inverting three have $3$-adic Kummer
completion
\[
                 \mu_3\oplus\mathbb Z_3\{1-\zeta_3\}.
\]
Set $u=\zeta_3(1-\zeta_3)$. The exact identities are
$u^2=-3$ and $c(u)=-u$. The Kummer class of $-1$
is zero at coefficient prime three, so the line of
$u$ is fixed by $c$; it is the free direct summand.
On the torsion summand $\mu_3$, conjugation acts by
inversion. The minus eigenspace is consequently
exactly $\mu_3$, and order-two descent identifies it
with $H^1(\mathcal G,\mathbb Z_3(1)\varepsilon_K)$.
This proves (E27.3).

These computations specify the coefficient group and
the bracket, quotient, and projector maps used to
separate the coordinates. The averaging denominator
$54$ is allowed only in the rational splitting
(E27.1); the integral torsion in E.24 and (E27.3)
is retained before that map. The anchor is
\S51.16, folios 733v--736v, scans 172--178,
and task LM7-E.23(i).
\end{proof}
